---
format: yohaku/0.2
target: text
macro: '[[macros/raw]]'
---

## 正文
````text
"""Loom-authoritative project discovery, planning and immutable source snapshots.

Projects are registered through the existing Workspace revision/expected contract.
The planner never executes project code or repairs generation by overwriting files.
The worker consumes only snapshots returned here, not a mutable world directory.
"""
import errno
import hashlib
import json
import os
from pathlib import Path, PurePosixPath
import shutil
import stat
import tempfile
import tomllib

from yohaku_project_schema import ProjectError, validate_manifest

PACKAGE = Path(__file__).resolve().parents[1]
REGISTRY = 'memory/projects.json'
REGISTRY_FORMAT = 'yohaku-project-registry/1'
SNAPSHOT_FORMAT = 'yohaku-project-snapshot/1'
MAX_FILE = 16 * 1024 * 1024
MAX_TOTAL = 64 * 1024 * 1024
RESERVED = {'.git', '.loom', '.yohaku', '.provenance'}


def digest(data):
    return hashlib.sha256(data).hexdigest()


def canonical(value):
    return json.dumps(value, ensure_ascii=False, sort_keys=True,
                      separators=(',', ':')).encode('utf-8')


def relative(name, *, metadata=False):
    if not isinstance(name, str) or not name or '\\' in name or '\0' in name:
        raise ProjectError('invalid_path', 'Use a nonempty POSIX relative path', path=name)
    p = PurePosixPath(name)
    # PurePath normalizes dots and duplicate separators; reject them before parsing.
    if p.is_absolute() or any(x in ('', '.', '..') for x in name.split('/')):
        raise ProjectError('invalid_path', 'Project paths must stay inside their root', path=name)
    if ':' in name or (not metadata and (p.name == '.project-snapshot.json' or any(x in RESERVED for x in p.parts))):
        raise ProjectError('invalid_path', 'Reserved metadata is not a project input', path=name)
    return p


def contained(root, name, *, metadata=False):
    p = relative(name, metadata=metadata)
    current = root
    for part in p.parts:
        current = current / part
        if current.is_symlink():
            raise ProjectError('symlink', 'Project paths cannot use symlinks', path=name)
    return current


def _read_bytes(root, name, *, metadata=False, limit=MAX_FILE):
    """Read ordinary bytes; Unix descriptors prevent ancestor/path substitution."""
    p = relative(name, metadata=metadata)
    if os.name == 'posix':
        fd = os.open(root, os.O_RDONLY | os.O_DIRECTORY | os.O_NOFOLLOW | os.O_CLOEXEC)
        try:
            for part in p.parts[:-1]:
                child = os.open(part, os.O_RDONLY | os.O_DIRECTORY | os.O_NOFOLLOW | os.O_CLOEXEC,
                                dir_fd=fd)
                os.close(fd)
                fd = child
            file_fd = os.open(p.name, os.O_RDONLY | os.O_NOFOLLOW | os.O_NONBLOCK | os.O_CLOEXEC, dir_fd=fd)
            try:
                before = os.fstat(file_fd)
                if not stat.S_ISREG(before.st_mode) or before.st_size > limit:
                    raise ProjectError('input_size', 'Input must be a bounded regular file', path=name)
                with os.fdopen(file_fd, 'rb', closefd=False) as stream:
                    data = stream.read(limit + 1)
                after = os.fstat(file_fd)
                if (before.st_dev, before.st_ino, before.st_size, before.st_mtime_ns) != (
                        after.st_dev, after.st_ino, after.st_size, after.st_mtime_ns):
                    raise ProjectError('source_conflict', 'Input changed while reading', path=name)
            finally:
                os.close(file_fd)
        finally:
            os.close(fd)
    else:
        # Windows source inspection is portable; this is not a native execution fallback.
        file = contained(root, name, metadata=metadata)
        before = file.stat()
        if not stat.S_ISREG(before.st_mode) or before.st_size > limit:
            raise ProjectError('input_size', 'Input must be a bounded regular file', path=name)
        with file.open('rb') as stream:
            data = stream.read(limit + 1)
        after = file.stat()
        contained(root, name, metadata=metadata)
        if (before.st_size, before.st_mtime_ns, before.st_ino) != (
                after.st_size, after.st_mtime_ns, after.st_ino):
            raise ProjectError('source_conflict', 'Input changed while reading', path=name)
    if len(data) > limit:
        raise ProjectError('input_size', 'Input exceeds project limit', path=name, limit=limit)
    return data


def read_bytes(root, name, *, metadata=False, limit=MAX_FILE):
    try:
        return _read_bytes(root, name, metadata=metadata, limit=limit)
    except OSError as error:
        kind = 'symlink' if error.errno == errno.ELOOP else 'missing_input' if isinstance(error, FileNotFoundError) else 'input_unavailable'
        raise ProjectError(kind, 'Project input could not be read', path=name, detail=str(error)) from error


class Projects:
    def __init__(self, world, workspace=None, state_dir=None):
        self.world = Path(world).resolve(strict=True)
        if workspace is None:
            from yohaku_agent import Workspace
            workspace = Workspace(self.world, PACKAGE / 'bin/loom-linux-x86_64')
        self.workspace = workspace
        self.state_dir = Path(state_dir) if state_dir is not None else (
            Path.home() / '.local/state/yohaku/project-snapshots' /
            digest(str(self.world).encode('utf-8')))

    def _execute(self, request):
        try:
            return self.workspace.execute(request)
        except ProjectError:
            raise
        except Exception as error:
            raise ProjectError('workspace_error', str(error), operation=request['command'],
                               replay_allowed=False) from error

    def _registry(self):
        file = contained(self.world, REGISTRY)
        if not file.exists():
            return {'format': REGISTRY_FORMAT, 'projects': []}, None
        captured = read_bytes(self.world, REGISTRY, limit=2 * 1024 * 1024)
        result = self._execute({'command': 'read_text', 'path': REGISTRY})
        text = result.get('text')
        if not isinstance(text, str) or text.encode('utf-8') != captured:
            raise ProjectError('source_conflict', 'Project registry changed during read')
        try:
            value = json.loads(text)
        except (ValueError, TypeError) as error:
            raise ProjectError('registry_invalid', 'Project registry must be JSON') from error
        if not isinstance(value, dict) or set(value) != {'format', 'projects'} or (
                value['format'] != REGISTRY_FORMAT or not isinstance(value['projects'], list)):
            raise ProjectError('registry_invalid', 'Unknown project registry format')
        seen = set()
        for row in value['projects']:
            if not isinstance(row, dict) or set(row) != {'id', 'source'} or (
                    not isinstance(row['id'], str) or row['id'] in seen):
                raise ProjectError('registry_invalid', 'Invalid or duplicate project registration')
            seen.add(row['id'])
            relative(row['source'])
        if len(seen) > 100:
            raise ProjectError('registry_invalid', 'Project registry exceeds 100 entries')
        return value, text

    def _entries(self):
        result = self._execute({'command': 'list'})
        files = result.get('files')
        if not isinstance(files, list):
            raise ProjectError('workspace_error', 'Workspace list omitted registered sources')
        return {row['source']: row for row in files if isinstance(row, dict) and isinstance(row.get('source'), str)}

    def _world_path(self, value, *, metadata=False):
        if not isinstance(value, str):
            raise ProjectError('generation_invalid', 'Loom metadata has no file path')
        # Normalize only trusted Workspace/Loom metadata, never caller manifest paths.
        # Rust canonicalize uses extended Windows paths; Python's world root may not.
        if os.name == 'nt':
            value = value.replace(chr(92), '/')
            if value.startswith('//?/UNC/'):
                value = '//' + value[8:]
            elif value.startswith('//?/'):
                value = value[4:]
        file = Path(value)
        if file.is_absolute():
            try:
                name = file.relative_to(self.world).as_posix()
            except ValueError as error:
                raise ProjectError('generation_invalid', 'Loom metadata points outside this world') from error
        else:
            name = value
        contained(self.world, name, metadata=metadata)
        return name

    def _document(self, source, entries, *, generation=True):
        relative(source)
        entry = entries.get(source)
        if not entry:
            raise ProjectError('unregistered_source', 'Register the authoritative tex with Workspace first',
                               source=source, recovery={'tool': 'adopt_file', 'path': source})
        captured_source = read_bytes(self.world, source, limit=2 * 1024 * 1024)
        document = self._execute({'command': 'read', 'source': source})
        if document.get('valid') is not True or not isinstance(document.get('text'), str):
            raise ProjectError('generation_invalid', 'Authoritative tex is invalid', source=source,
                               detail=document.get('error'))
        source_data = document['text'].encode('utf-8')
        if captured_source != source_data or read_bytes(self.world, source) != source_data:
            raise ProjectError('source_conflict', 'Source changed after Workspace read', source=source)
        if not generation:
            return {'source': source, 'source_data': source_data, 'source_hash': digest(source_data),
                    'body': document.get('body')}
        synced = entry.get('last_sync') or {}
        if not isinstance(synced, dict) or not isinstance(synced.get('result', {}), dict):
            raise ProjectError('generation_invalid', 'Invalid Workspace generation receipt', source=source)
        build = synced.get('result', {}).get('build') or {}
        if not isinstance(build, dict):
            raise ProjectError('generation_invalid', 'Invalid Loom build receipt', source=source)
        if synced.get('ok') is not True or synced.get('result', {}).get('state') != 'generated':
            raise ProjectError('generation_required', 'Source has no successful current Loom generation',
                               source=source, generation=synced,
                               recovery={'tool': 'sync_document', 'source': source})
        try:
            map_name = self._world_path(build.get('map_path'), metadata=True)
            map_bytes = read_bytes(self.world, map_name, metadata=True)
            mapping = json.loads(map_bytes)
        except ProjectError:
            raise
        except (OSError, ValueError) as error:
            raise ProjectError('generation_invalid', 'Loom source map is unavailable', source=source) from error
        if not isinstance(mapping, dict) or mapping.get('format') != 'loom-source-map/0.1':
            raise ProjectError('generation_invalid', 'Unsupported Loom source map', source=source)
        output = self._world_path(entry.get('output'))
        if self._world_path(mapping.get('source')) != source or self._world_path(mapping.get('output')) != output:
            raise ProjectError('generation_invalid', 'Loom map and Workspace routing differ', source=source)
        macro = self._world_path(mapping.get('macro'))
        macro_data = read_bytes(self.world, macro)
        output_data = read_bytes(self.world, output)
        current = {'source_hash': digest(source_data), 'macro_hash': digest(macro_data),
                   'output_hash': digest(output_data)}
        changes = {key: {'expected': mapping.get(key), 'actual': value}
                   for key, value in current.items() if mapping.get(key) != value}
        if changes:
            raise ProjectError('generation_drift', 'Authoritative input or generated output differs from its Loom build',
                               source=source, output=output, changes=changes, preserved=True,
                               recovery={'tool': 'read_document', 'source': source})
        if mapping.get('source_snapshot') != document['text'] or (
                not isinstance(mapping.get('macro_snapshot'), str) or
                mapping['macro_snapshot'].encode('utf-8') != macro_data):
            raise ProjectError('generation_invalid', 'Loom map snapshots do not match its current input hashes', source=source)
        return {'source': source, 'output': output, 'macro': macro,
                'source_data': source_data, 'output_data': output_data, 'macro_data': macro_data,
                'map_data': map_bytes, 'map': mapping, 'body': document.get('body'), **current}

    def _manifest(self, source, entries=None, *, generation=True):
        entries = self._entries() if entries is None else entries
        doc = self._document(source, entries, generation=generation)
        try:
            value = json.loads(doc['body'])
        except (ValueError, TypeError) as error:
            raise ProjectError('manifest_invalid', 'Project manifest body must be JSON') from error
        manifest = validate_manifest(value)
        prefix = manifest['root'].rstrip('/') + '/'
        if not source.startswith(prefix):
            raise ProjectError('manifest_invalid', 'Manifest source must be inside its declared project root', source=source)
        return manifest, doc

    def register(self, manifest_source):
        relative(manifest_source)
        if not manifest_source.endswith('.tex'):
            raise ProjectError('manifest_invalid', 'Register the authoritative manifest .tex, not generated JSON')
        manifest, doc = self._manifest(manifest_source)
        registry, expected = self._registry()
        old = next((row for row in registry['projects'] if row['id'] == manifest['id']), None)
        if old and old['source'] != manifest_source:
            raise ProjectError('project_conflict', 'Project id is already bound to another manifest', project_id=manifest['id'])
        if any(row['source'] == manifest_source and row['id'] != manifest['id'] for row in registry['projects']):
            raise ProjectError('project_conflict', 'Manifest is registered under another project id')
        if old is None:
            registry['projects'].append({'id': manifest['id'], 'source': manifest_source})
            text = json.dumps(registry, ensure_ascii=False, indent=2) + '\n'
            self._execute({'command': 'write_text', 'path': REGISTRY, 'expected': expected, 'text': text})
            actual, _ = self._registry()
            if {'id': manifest['id'], 'source': manifest_source} not in actual['projects']:
                raise ProjectError('outcome_unknown', 'Registration read-back mismatch; query before retry', replay_allowed=False)
        return {'project_id': manifest['id'], 'manifest_source': manifest_source,
                'manifest_revision': doc['source_hash'], 'registered': True, 'duplicate': old is not None,
                'executed': False}

    def list(self):
        registry, _ = self._registry()
        rows = []
        for row in registry['projects']:
            try:
                info = self.describe(row['id'])
                rows.append({'project_id': row['id'], 'manifest_source': row['source'],
                             'root': info['project']['root'], 'runtime': info['project']['runtime'],
                             'tasks': info['tasks'], 'available': info['available'],
                             'generation': info['generation']})
            except ProjectError as error:
                rows.append({'project_id': row['id'], 'manifest_source': row['source'],
                             'available': False, 'error': error.to_dict()})
        return {'format': REGISTRY_FORMAT, 'projects': rows}

    def _description(self, project_id):
        registry, _ = self._registry()
        row = next((row for row in registry['projects'] if row['id'] == project_id), None)
        if row is None:
            raise ProjectError('project_not_found', 'Unknown project_id; register a manifest first', project_id=project_id)
        manifest, doc = self._manifest(row['source'], generation=False)
        if manifest['id'] != project_id:
            raise ProjectError('project_conflict', 'Registered manifest changed its project id', project_id=project_id)
        return {'format': 'yohaku-project-description/1', 'project_id': project_id,
                'project': manifest, 'manifest': manifest, 'manifest_source': row['source'],
                'manifest_revision': doc['source_hash'], 'revision': doc['source_hash'], 'tasks': sorted(manifest['tasks']),
                'editable_sources': [manifest['root'] + '/' + source for source in manifest['sources']],
                'output_root': '/output', 'executed': False}

    def describe(self, project_id):
        # Discovery stays useful when a source needs repair. Only plan/materialize
        # require a clean generation barrier; no bad input is allowed to execute.
        info = self._description(project_id)
        try:
            current, _, mappings, source_set = self._collect(project_id)
            return {**current, 'source_set': source_set, 'mappings': mappings,
                    'available': True, 'generation': {'ok': True}}
        except ProjectError as error:
            return {**info, 'source_set': None, 'mappings': [], 'available': False,
                    'generation': error.to_dict()}

    def _collect(self, project_id, expected_sources=None):
        info = self._description(project_id)
        project = info['project']
        entries = self._entries()
        prefix = project['root'].rstrip('/') + '/'
        names = [info['manifest_source']] + [prefix + source for source in project['sources']]
        names = list(dict.fromkeys(names))
        files = {}
        mappings = []
        checks = []
        total = 0

        def put(name, data, role, *, internal=False, **fields):
            nonlocal total
            relative(name, metadata=internal)
            if any(name.startswith(existing + '/') or existing.startswith(name + '/') for existing in files):
                raise ProjectError('input_conflict', 'Project file and directory paths overlap', path=name)
            if name in files:
                if files[name]['data'] != data or files[name]['role'] != role:
                    raise ProjectError('input_conflict', 'Two project inputs have the same snapshot path', path=name)
                return
            total += len(data)
            if total > MAX_TOTAL:
                raise ProjectError('input_size', 'Project snapshot exceeds 64 MiB', limit=MAX_TOTAL)
            files[name] = {'path': name, 'data': data, 'sha256': digest(data),
                           'bytes': len(data), 'role': role, **fields}

        for source in names:
            doc = self._document(source, entries)
            if not doc['output'].startswith(prefix):
                raise ProjectError('output_outside_project', 'Declared source generates outside its project root',
                                   source=source, output=doc['output'])
            source_rel = source[len(prefix):]
            output_rel = doc['output'][len(prefix):]
            put(source_rel, doc['source_data'], 'source', source=source)
            put(output_rel, doc['output_data'], 'generated', source=source,
                source_revision=doc['source_hash'])
            map_name = '.provenance/maps/' + digest(source.encode()) + '.json'
            # Only internally constructed provenance names may enter the reserved subtree.
            put(map_name, doc['map_data'], 'map', internal=True)
            macro_name = '.provenance/macros/' + doc['macro']
            put(macro_name, doc['macro_data'], 'macro', internal=True)
            mappings.append({'source': source, 'source_revision': doc['source_hash'],
                             'output': output_rel, 'output_hash': doc['output_hash'],
                             'macro': doc['macro'], 'macro_hash': doc['macro_hash'],
                             'map': map_name,
                             'generated_to_source_line': doc['map'].get('generated_to_source_line'),
                             'body_source_line': doc['map'].get('body_source_line')})
            checks.extend([(source, doc['source_hash']), (doc['output'], doc['output_hash']),
                           (doc['macro'], doc['macro_hash'])])
        for asset in project.get('assets', []):
            name = asset['path'] if isinstance(asset, dict) else asset
            data = read_bytes(self.world, prefix + name)
            expected = asset.get('sha256') if isinstance(asset, dict) else None
            if expected is not None and digest(data) != expected:
                raise ProjectError('asset_drift', 'Asset does not match its declared hash', path=name,
                                   expected=expected, actual=digest(data), preserved=True)
            put(name, data, 'asset')
            checks.append((prefix + name, digest(data)))
        # Fingerprint all captured bytes and manifest revision, not just source mtimes.
        fingerprint = {'project_id': project_id, 'manifest_revision': info['manifest_revision'],
                       'files': [{key: row[key] for key in ('path', 'sha256', 'bytes', 'role')}
                                 for row in sorted(files.values(), key=lambda x: x['path'])]}
        source_set = digest(canonical(fingerprint))
        if expected_sources is not None and expected_sources != source_set:
            raise ProjectError('source_conflict', 'Project sources changed since planning',
                               expected_sources=expected_sources, actual_sources=source_set,
                               recovery={'tool': 'project_plan', 'project_id': project_id})
        for name, expected in checks:
            if digest(read_bytes(self.world, name)) != expected:
                raise ProjectError('source_conflict', 'Project changed while collecting inputs', path=name)
        return info, files, mappings, source_set

    def _runtime(self, project):
        missing = []
        runtime = {'kind': project['runtime'], 'platform': 'linux', 'offline': True}
        if os.name != 'posix' or not os.access('/usr/bin/bwrap', os.X_OK):
            missing.append({'kind': 'runtime_unavailable', 'resource': 'linux-bubblewrap',
                            'message': 'Linux execution requires the existing mandatory bwrap backend; no native fallback'})
        if project['runtime'] == 'python':
            runtime['program'] = '/usr/bin/python3'
            if os.name != 'posix' or not os.access('/usr/bin/python3', os.X_OK):
                missing.append({'kind': 'runtime_missing', 'resource': '/usr/bin/python3'})
        else:
            try:
                from yohaku_dev import validate_profile
                profile = json.loads((PACKAGE / 'config/development-runtime.json').read_text('utf-8'))
                toolchain, registry = validate_profile(profile)
                runtime.update(toolchain=str(toolchain), registry=str(registry),
                               program='/opt/yohaku/rust/bin/cargo',
                               fingerprint=digest(canonical(profile)))
            except Exception as error:
                missing.append({'kind': 'runtime_missing', 'resource': 'rust-offline-profile', 'message': str(error)})
        return runtime, missing

    def _cargo_dependencies(self, files, runtime, *, details=False):
        """Cache evidence, not a Cargo cfg/features/semver resolver.

        A unique locked version of a definitely-required direct registry dependency
        may fail preflight when missing. Other missing lock entries are unresolved
        warnings: only the mandatory isolated Cargo --offline --locked run determines
        activation. A present corrupt archive is always fatal, even if inactive.
        """
        report = {'missing': [], 'dependency_warnings': [], 'cache_coverage_partial': False,
                  'required_direct_packages': [], 'coverage': {
                      'activation_resolved': False, 'registry_packages': 0, 'cached_packages': 0,
                      'scope': 'top-level non-optional dependencies/build-dependencies/dev-dependencies; aliases use package name; no cfg/features/semver resolution',
                      'final_resolver': 'mandatory isolated Cargo --offline --locked'}}
        missing = report['missing']
        warnings = report['dependency_warnings']
        def finish():
            report['cache_coverage_partial'] = bool(warnings)
            return report if details else missing
        if 'Cargo.toml' not in files or 'Cargo.lock' not in files:
            missing.append({'kind': 'dependency_manifest_missing', 'resource': 'Cargo.toml and Cargo.lock',
                            'message': 'Declare their authoritative tex sources for an offline locked build'})
            return finish()
        try:
            manifest = tomllib.loads(files['Cargo.toml']['data'].decode('utf-8'))
            lock = tomllib.loads(files['Cargo.lock']['data'].decode('utf-8'))
        except (ValueError, UnicodeError) as error:
            missing.append({'kind': 'dependency_manifest_invalid', 'message': str(error)})
            return finish()
        required = set()
        tables = []
        groups = ('dependencies', 'dev-dependencies', 'build-dependencies')
        for group in groups:
            table = manifest.get(group, {})
            if not isinstance(table, dict):
                missing.append({'kind': 'dependency_manifest_invalid', 'resource': group})
                return finish()
            tables.append((group, table, True))
        targets = manifest.get('target', {})
        if not isinstance(targets, dict):
            missing.append({'kind': 'dependency_manifest_invalid', 'resource': 'target'})
            return finish()
        for target, target_tables in targets.items():
            if not isinstance(target_tables, dict):
                missing.append({'kind': 'dependency_manifest_invalid', 'resource': 'target.' + str(target)})
                return finish()
            for group in groups:
                table = target_tables.get(group, {})
                if not isinstance(table, dict):
                    missing.append({'kind': 'dependency_manifest_invalid', 'resource': 'target.' + str(target) + '.' + group})
                    return finish()
                tables.append(('target.' + str(target) + '.' + group, table, False))
        # Patches can change a registry dependency into local/alternate input. Do
        # not claim its cache version is definitely required without resolving it.
        patched = set()
        patches = manifest.get('patch', {})
        if not isinstance(patches, dict):
            missing.append({'kind': 'dependency_manifest_invalid', 'resource': 'patch'})
            return finish()
        for table in patches.values():
            if not isinstance(table, dict):
                missing.append({'kind': 'dependency_manifest_invalid', 'resource': 'patch table'})
                return finish()
            for alias, value in table.items():
                patched_name = value.get('package', alias) if isinstance(value, dict) else alias
                if not isinstance(patched_name, str) or not patched_name:
                    missing.append({'kind': 'dependency_manifest_invalid', 'resource': 'patch package name'})
                    return finish()
                patched.add(patched_name)
                tables.append(('patch', {alias: value}, False))
        for group, table, top_level in tables:
            for alias, value in table.items():
                if isinstance(value, str):
                    if not value:
                        missing.append({'kind': 'dependency_manifest_invalid', 'resource': group + '.' + alias + '.version'})
                        continue
                    name, optional, is_registry = alias, False, True
                elif isinstance(value, dict):
                    name = value.get('package', alias)
                    optional = value.get('optional', False)
                    if 'version' in value and (not isinstance(value['version'], str) or not value['version']):
                        missing.append({'kind': 'dependency_manifest_invalid', 'resource': group + '.' + alias + '.version'})
                        continue
                    if type(optional) is not bool:
                        missing.append({'kind': 'dependency_manifest_invalid', 'resource': group + '.' + alias + '.optional'})
                        continue
                    is_registry = not any(key in value for key in ('path', 'git', 'workspace'))
                    if 'path' in value:
                        try:
                            dependency_path = relative(value['path'])
                            candidate = (dependency_path / 'Cargo.toml').as_posix()
                            if candidate not in files:
                                item = {'kind': 'dependency_missing', 'resource': candidate,
                                        'message': 'Declared path dependency is not in the fixed input snapshot'}
                                (missing if top_level and not optional else warnings).append(item)
                        except ProjectError:
                            missing.append({'kind': 'dependency_path_outside', 'resource': value['path']})
                else:
                    missing.append({'kind': 'dependency_manifest_invalid', 'resource': group + '.' + alias})
                    continue
                if not isinstance(name, str) or not name or any(x in name for x in ('/', '\\', ':', '\0')):
                    missing.append({'kind': 'dependency_manifest_invalid', 'resource': group + '.' + alias + '.package'})
                    continue
                if is_registry and top_level and not optional and name not in patched:
                    required.add(name)
        report['required_direct_packages'] = sorted(required)
        packages = lock.get('package', [])
        if not isinstance(packages, list) or any(not isinstance(row, dict) for row in packages):
            missing.append({'kind': 'dependency_manifest_invalid', 'resource': 'Cargo.lock package array'})
            return finish()
        registry_packages = []
        package_ids = set()
        for package in packages:
            name, version, source = package.get('name'), package.get('version'), package.get('source')
            if not isinstance(name, str) or not isinstance(version, str) or (
                    not name or not version or any(x in name + version for x in ('/', '\\', ':', '\0'))):
                missing.append({'kind': 'dependency_manifest_invalid', 'resource': 'package name/version'})
                continue
            package_id = (name, version, source) if isinstance(source, (str, type(None))) else None
            if package_id is None or package_id in package_ids:
                missing.append({'kind': 'dependency_manifest_invalid', 'resource': 'duplicate or invalid package identity'})
                continue
            package_ids.add(package_id)
            if source is None:
                continue
            if not isinstance(source, str) or not source.startswith('registry+'):
                missing.append({'kind': 'dependency_source_unsupported', 'resource': source})
                continue
            checksum = package.get('checksum')
            if not isinstance(checksum, str) or len(checksum) != 64 or any(c not in '0123456789abcdef' for c in checksum):
                missing.append({'kind': 'dependency_manifest_invalid', 'resource': name + '@' + version + ' checksum'})
                continue
            registry_packages.append(package)
        report['coverage']['registry_packages'] = len(registry_packages)
        candidates = {}
        for package in registry_packages:
            candidates.setdefault(package['name'], []).append(package)
        for name in required:
            if name not in candidates:
                missing.append({'kind': 'dependency_not_locked', 'resource': name,
                                'message': 'Definitely-required direct registry dependency is absent from Cargo.lock'})
        if 'registry' not in runtime:
            warnings.append({'kind': 'dependency_cache_unresolved', 'message': 'Prepared registry unavailable; runtime preflight reports the fatal runtime requirement'})
            return finish()
        registry = Path(runtime['registry'])
        try:
            cache_dirs = [p for p in (registry / 'cache').iterdir() if p.is_dir()]
            source_dirs = [p for p in (registry / 'src').iterdir() if p.is_dir()]
        except OSError as error:
            missing.append({'kind': 'dependency_cache_unavailable', 'message': type(error).__name__})
            return finish()
        for package in registry_packages:
            name, version, checksum = package['name'], package['version'], package['checksum']
            archives = [p / (name + '-' + version + '.crate') for p in cache_dirs
                        if (p / (name + '-' + version + '.crate')).is_file()]
            unpacked = next((p / (name + '-' + version) for p in source_dirs
                             if (p / (name + '-' + version) / '.cargo-checksum.json').is_file()), None)
            available = False
            corrupt = False
            for archive in archives:
                try:
                    with archive.open('rb') as stream:
                        actual = hashlib.file_digest(stream, 'sha256').hexdigest()
                except OSError as error:
                    missing.append({'kind': 'dependency_cache_unavailable', 'resource': name + '@' + version,
                                    'message': type(error).__name__})
                    corrupt = True
                    continue
                if actual != checksum:
                    missing.append({'kind': 'dependency_checksum_mismatch', 'resource': name + '@' + version,
                                    'expected_sha256': checksum, 'actual_sha256': actual,
                                    'message': 'Existing archive checksum differs from Cargo.lock; no fallback to unpacked source'})
                    corrupt = True
                else:
                    available = True
            if corrupt:
                continue
            if not available and unpacked is not None:
                try:
                    source_checksum = json.loads((unpacked / '.cargo-checksum.json').read_text('utf-8'))
                    available = source_checksum.get('package') == checksum
                except (OSError, ValueError, AttributeError):
                    available = False
            if available:
                report['coverage']['cached_packages'] += 1
                continue
            item = {'kind': 'dependency_missing' if name in required and len(candidates[name]) == 1 else 'dependency_cache_unresolved',
                    'resource': name + '@' + version,
                    'message': 'Not available in the prepared read-only offline Cargo cache'}
            if name in required and len(candidates[name]) == 1:
                missing.append(item)
            else:
                item.update(activation='unresolved',
                            reason='ambiguous locked versions' if name in required else 'optional, target-specific or lockfile-only activation is not resolved',
                            final_resolver='mandatory isolated Cargo --offline --locked')
                warnings.append(item)
        return finish()

    def plan(self, project_id, action, expected_sources=None):
        info, files, mappings, source_set = self._collect(project_id, expected_sources)
        project = info['project']
        task = project['tasks'].get(action)
        if task is None:
            raise ProjectError('task_not_found', 'Unknown project action', project_id=project_id,
                               action=action, available=sorted(project['tasks']))
        runtime, missing = self._runtime(project)
        args = task.get('args', [])
        result = {'format': 'yohaku-project-plan/1', 'project_id': project_id, 'action': action,
                  'project': project, 'manifest': project, 'manifest_revision': info['manifest_revision'],
                  'source_set': source_set, 'runtime': runtime, 'cwd': '/project',
                  'read_only': ['/project'], 'writable': ['/output'],
                  'capabilities': {'network': 'isolated', 'filesystem': 'declared-snapshot-only'},
                  'artifacts': task.get('artifacts', []), 'missing': missing,
                  'executed': False, 'mappings': mappings, 'dependency_warnings': [],
                  'cache_coverage_partial': False}
        if task['kind'] == 'python':
            if task['entry'] not in files or files[task['entry']]['role'] != 'generated':
                missing.append({'kind': 'entry_missing', 'resource': task['entry'],
                                'message': 'Python entry must be a successfully generated declared source'})
            result['argv'] = ['/usr/bin/python3', '-B', '/project/' + task['entry'], *args]
        elif task['kind'] == 'cargo':
            dependencies = self._cargo_dependencies(files, runtime, details=True)
            missing.extend(dependencies['missing'])
            result['dependency_warnings'] = dependencies['dependency_warnings']
            result['cache_coverage_partial'] = dependencies['cache_coverage_partial']
            result['dependency_coverage'] = dependencies['coverage']
            managed_args = list(args)
            for flag in ('--offline', '--locked'):
                if flag not in managed_args:
                    managed_args.append(flag)
            result['argv'] = ['/opt/yohaku/rust/bin/cargo', *managed_args]
            result['dependency_preflight'] = 'Partial cache evidence: definitely-required direct non-optional package versions may fail preflight; optional/target/transitive activation and ambiguous versions are unresolved warnings. Existing corrupt archives remain fatal. No cfg/features/semver resolver; only the mandatory isolated Cargo --offline --locked run resolves and verifies the active dependency graph.'
        else:
            result['argv'] = ['/output/' + task['artifact'], *args]
            result['requires_artifact'] = {'action': task['from'], 'path': task['artifact'],
                                           'source_set': source_set}
        result['ok'] = not missing
        return result

    def materialize(self, project_id, expected_sources=None):
        info, files, mappings, source_set = self._collect(project_id, expected_sources)
        self.state_dir.mkdir(mode=0o700, parents=True, exist_ok=True)
        if self.state_dir.is_symlink() or not self.state_dir.is_dir():
            raise ProjectError('state_invalid', 'Snapshot state directory must be an ordinary private directory')
        if os.name == 'posix':
            meta = self.state_dir.stat()
            if meta.st_uid != os.getuid() or meta.st_mode & 0o077:
                raise ProjectError('state_invalid', 'Snapshot state must be owned and private')
        target = self.state_dir / source_set
        manifest = {'format': SNAPSHOT_FORMAT, 'id': source_set, 'source_set': source_set,
                    'project_id': project_id, 'manifest_revision': info['manifest_revision'],
                    'files': [{key: value for key, value in row.items() if key != 'data'}
                              for row in sorted(files.values(), key=lambda x: x['path'])],
                    'mappings': mappings}
        if target.is_symlink():
            raise ProjectError('state_invalid', 'Snapshot identity is a symlink')
        if target.exists():
            try:
                recorded = json.loads(read_bytes(target, '.project-snapshot.json', metadata=True))
            except (OSError, ValueError) as error:
                raise ProjectError('snapshot_conflict', 'Existing snapshot has no valid receipt') from error
            if recorded != manifest:
                raise ProjectError('snapshot_conflict', 'Existing snapshot receipt differs from source set')
            expected_files = {row['path'] for row in manifest['files']} | {'.project-snapshot.json'}
            actual_files = set()
            for folder, dirs, names in os.walk(target, followlinks=False):
                if any((Path(folder) / name).is_symlink() for name in dirs + names):
                    raise ProjectError('snapshot_conflict', 'Existing snapshot contains a symlink')
                actual_files.update((Path(folder) / name).relative_to(target).as_posix() for name in names)
            if actual_files != expected_files:
                raise ProjectError('snapshot_conflict', 'Existing snapshot file set was modified')
            for row in manifest['files']:
                if digest(read_bytes(target, row['path'], metadata=True)) != row['sha256']:
                    raise ProjectError('snapshot_conflict', 'Existing immutable snapshot was modified', path=row['path'])
            return {**manifest, 'path': str(target), 'reused': True}
        temporary = Path(tempfile.mkdtemp(prefix='.preparing-', dir=self.state_dir))
        try:
            for row in files.values():
                dest = contained(temporary, row['path'], metadata=True)
                dest.parent.mkdir(parents=True, exist_ok=True)
                with dest.open('xb') as stream:
                    stream.write(row['data'])
                    stream.flush()
                    os.fsync(stream.fileno())
                if digest(dest.read_bytes()) != row['sha256']:
                    raise ProjectError('snapshot_conflict', 'Snapshot copy verification failed', path=row['path'])
                if os.name == 'posix':
                    dest.chmod(0o444)
            receipt = temporary / '.project-snapshot.json'
            with receipt.open('xb') as stream:
                stream.write(canonical(manifest) + b'\n')
                stream.flush()
                os.fsync(stream.fileno())
            # Recheck the current whole source set after the copy, not just mtimes.
            self._collect(project_id, source_set)
            try:
                temporary.rename(target)
                if os.name == 'posix':
                    descriptor = os.open(self.state_dir, os.O_RDONLY | os.O_DIRECTORY | os.O_CLOEXEC)
                    try:
                        os.fsync(descriptor)
                    finally:
                        os.close(descriptor)
            except OSError as error:
                if error.errno in (errno.EEXIST, errno.ENOTEMPTY) and target.exists():
                    return self.materialize(project_id, source_set)
                raise
        finally:
            if temporary.exists():
                shutil.rmtree(temporary)
        return {**manifest, 'path': str(target), 'reused': False}
````

## 输出与关联
```yaml
output: "research/workflow-runtime/agent/yohaku_project_sources.py"
```
